\begin{mistake}{2026-05-27}{06}

\begin{problemcontent}
设 \(a>0\)，\(f(x)=g(x)=\begin{cases} a, & 0 \leqslant x \leqslant 1, \\ 0, & \text{其他.} \end{cases}\) \(D\) 表示全平面，则 \(\iint_D f(x)g(y-x) \mathrm{d}x\mathrm{d}y = \underline{\quad\quad}\)
\end{problemcontent}

\begin{solutioncontent}
\textbf{解法一：确定有效积分区域}

被积函数不为 0 的条件为：
\(\begin{cases} 0 \leqslant x \leqslant 1 \\ 0 \leqslant y-x \leqslant 1 \end{cases} \Rightarrow \begin{cases} 0 \leqslant x \leqslant 1 \\ x \leqslant y \leqslant x+1 \end{cases}\)

在此区域 \(\Omega\) 内，被积函数为 \(a \cdot a = a^2\)，转化为累次积分：
\[\iint_D f(x)g(y-x) \mathrm{d}x\mathrm{d}y = \int_0^1 \mathrm{d}x \int_x^{x+1} a^2 \mathrm{d}y = \int_0^1 a^2 \mathrm{d}x = a^2\]

\textbf{解法二：平移换元与变量分离}

化为累次积分：
\[\iint_D f(x)g(y-x) \mathrm{d}x\mathrm{d}y = \int_{-\infty}^{+\infty} f(x) \left( \int_{-\infty}^{+\infty} g(y-x) \mathrm{d}y \right) \mathrm{d}x\]
内层积分令 \(u = y-x\)，则 \(\mathrm{d}u = \mathrm{d}y\)：
\[\int_{-\infty}^{+\infty} g(y-x) \mathrm{d}y = \int_{-\infty}^{+\infty} g(u) \mathrm{d}u = \int_0^1 a \mathrm{d}u = a\]
原式 \(= \left(\int_0^1 a \mathrm{d}x\right) \cdot a = a \cdot a = a^2\)
\end{solutioncontent}

\mistakeknowledge{
\begin{itemize}
  \item \textbf{核心公式/定理}：二重积分化为累次积分（Fubini定理）；一元积分平移换元。
  \item \textbf{易错点}：寻找非零区域时，误将 \(0 \leqslant y-x \leqslant 1\) 直接作为 \(y\) 的独立界限，导致区域形状错误。
  \item \textbf{关联知识}：概率论中连续型随机变量和 \(Z=X+Y\) 的概率密度计算（卷积公式）。
\end{itemize}}

\end{mistake}
